%%=============================================================================
%% Analyse en ontwerp
%%=============================================================================

\chapter{\IfLanguageName{dutch}{Analyse en ontwerp}{Analysis and design}}%
\label{ch:ontwerp}

%% Richtlengte: 2 tot 3 bladzijden.
%%
%% Hier toon je dat er nagedacht is voordat er getypt werd. Dit hoofdstuk
%% weegt zwaar: het is waar het verschil zichtbaar wordt tussen een uitgevoerde
%% taak en een zelfstandig project.

Vervang deze tekst door een korte inleiding op het hoofdstuk.

\section{\IfLanguageName{dutch}{Vereisten}{Requirements}}%
\label{sec:vereisten}

%% Wat moest het worden? Onderscheid wat noodzakelijk was van wat wenselijk
%% was, en vermeld bij wie je die vereisten opgehaald hebt (je mentor, het
%% team, de gebruikers).
%%
%% Vermeld ook de niet-functionele vereisten die er echt toe deden:
%% snelheid, veiligheid, privacy, onderhoudbaarheid, beschikbaarheid.

\section{\IfLanguageName{dutch}{Architectuur}{Architecture}}%
\label{sec:architectuur}

%% Toon de opbouw van je oplossing met een figuur: welke onderdelen zijn er,
%% wat doet elk onderdeel, en hoe praten ze met elkaar? Een goed getekend
%% schema vervangt een bladzijde tekst.
%%
%% Plaats je afbeeldingen in de map graphics/.

\begin{figure}[htbp]
  \centering
  \includegraphics[width=0.8\textwidth]{voorbeeldafbeelding}
  \caption[Architectuur van de oplossing]{\label{fig:architectuur}%
    Vervang deze afbeelding door je eigen architectuurschema. Zorg voor een
    bijschrift dat de figuur volledig beschrijft, zodat een lezer ze begrijpt
    zonder in de tekst te moeten zoeken. Neem je een figuur over van iemand
    anders, vermeld dan de bron.}
\end{figure}

\section{\IfLanguageName{dutch}{Gegevensmodel}{Data model}}%
\label{sec:gegevensmodel}

%% Werk je met gegevens, toon dan hoe ze gestructureerd zijn en waarom.
%% Heeft je project geen noemenswaardig gegevensmodel, verwijder deze sectie
%% dan gerust.

\section{\IfLanguageName{dutch}{Ontwerpkeuzes}{Design decisions}}%
\label{sec:ontwerpkeuzes}

%% Noteer drie tot vijf keuzes die je bewust gemaakt hebt, telkens in deze
%% vorm: de keuze, de alternatieven, de reden, en het gevolg.
%%
%% Dit is de sectie waaruit een jurylid zijn vragen haalt. Een keuze waarvan
%% je het alternatief niet kan benoemen, was geen keuze.
